\documentclass[10pt]{article}

\usepackage[utf8]{inputenc}

\usepackage[T1]{fontenc}

\usepackage{amsmath}

\usepackage{amsfonts}

\usepackage{amssymb}

\usepackage[version=4]{mhchem}

\usepackage{stmaryrd}

\usepackage{bbold}



\begin{document}

стве стандартного слоя пространство $T^{p, q}\left(\mathbb{R}^{n}\right)$ тензоров типа $(p, q)$ на $\mathbb{R} n$, в к-ром группа $\mathrm{G} \mathrm{L}(n, \mathbb{R})$ действует при помощи тензорного представления. Напр., $T^{1,0}(M)$ совпадает с касательным расслоением $T(M)$ над $M$, а $T^{0,1}(M)$ - с кокасательным расслоением $T(M)^{*}$. В общем случае Т. р. изоморфно тензорному произведению касательных и кокасательных расслоений:



\[

T^{p, q}(M) \cong \bigotimes^{p} T(M) \otimes \otimes^{q} T(M)^{*} .

\]



Сечения Т. $\mathbf{p}$. типа $(p, q)$ наз. т е н з о р н ы м и п ол я м и типа $(p, q)$ и являются основным объектом исследования дифференциальной геометрии. Напр., риманова структура на $M$ - это гладкое сечение расслоения $T^{0,2}(M)$, значения к-рого являются положительно определенными симметрич. формами. Гладкие сечения расслоения $T^{p, q}(M)$ образуют модуль $D^{p, q}(M)$ над алгеброй $F^{\infty}(M)=D^{0,0}(M)$ гладких функций на $M$. Если $M$ - паракомпактное хаусдорфово многообразие, то



\[

D^{p, q}(M) \cong \stackrel{p}{\otimes} D^{1}(M) \otimes \stackrel{q}{\bigotimes} D^{1}(M)^{*},

\]



где $D^{1}(M)=D^{1,0}(M)$ - модуль гладких векторных полей, $D^{1}(M)^{*}=D^{0,1}(M)$ - модуль пфаффовых дифференциальных форм, а тензорные произведения берутся над $F^{\infty}(M)$. В классической дифференциальной геометрии тензорные поля иногда наз. просто т е н 3 ор а м и на $M$.



Лит.: [1] К о б а я с и Ші., Н о м и д з у К., Основы дифференциальной геометрии, т. 1, пер. с англ., М., 1981; [2] Х е лг а с о н С., Дифференциальная геометрия и симметрические

пространства, пер. с англ. М., 1964.



ТЕНЗОРНЫЙ АНАЛИЗ - обобщение векторного анализа, раздел тензорного исчисления, изучающий дифференциальные операторы, действующие на алгебре тензорных полей $D(M)$ диффференцируемого многообразия $M$. Рассматриваются также операторы, действующие на более общие, чем тензорные поля, геометрич. объекты: тензорные плотности, дифференциальные формы со значениями в векторном расслоении у т. Д.



Наибольший интерес представляют операторы, действие к-рых не выводит за пределы алгебры $D(M)$.



\begin{enumerate}

  \item Ковариантная прочзводная вдоль векторного поля $X$ - линейное отображение $\nabla_{X}$ пространства векторных полей $D^{1}(M)$ многообразия $M$, зависящее от векторного поля $X$ и удовлетворяющее условиям:

\end{enumerate}



$\nabla_{f} X+g Y^{Z}=f \nabla X^{Z}+g \nabla_{Y} Z, \quad \nabla X(f Z)=f \nabla X_{X} Z+(X f) Y$, где $X, Y, Z \in D^{\prime}(M), f, g$ - гладкие функции на $M$. Определяемые этим оператором связность Г и параллельное перенесение позволяют распространить действие ковариантной производной до линейного отображения алгебры $D(M)$ в себя; при этом отображение $\nabla X$ есть дифференцирование, сохраняет тиц тензорного поля и перестановочно со сверткой.



В локальных координатах $u^{1}, u^{2}, \ldots, u^{n}$ ковариантная производная тензора с компонентами $T\left(T_{j_{1} \ldots j_{m}}^{i_{1} \ldots i_{l}}\right)$ относительно вектора $X=\xi i \frac{\partial}{\partial u^{i}}$ определяется так:



\[

\begin{gathered}

\nabla_{X} T=\xi s\left(\frac{\partial T_{j_{1} \ldots j_{m}}^{i_{1} \ldots i_{l}}}{\partial u^{s}}+\Gamma_{k s}^{i_{1}} T_{i_{1} \ldots i_{m}}^{k \ldots i_{l}}+\ldots\right. \\

\left.\ldots-\Gamma_{j_{1}, s}^{k} T_{k \ldots i_{m}}^{i_{1} \ldots i_{l}}-\ldots\right)

\end{gathered}

\]



$\Gamma_{k s}^{i}$ - объект связности $\Gamma$.



\begin{enumerate}

  \setcounter{enumi}{1}

  \item Ли производная вдоль векторного поля $X$-отображение $L_{X}$ пространства $D^{\prime}(M)$, определяемое формулой $L_{X}: Y \rightarrow[X, Y]$, где $[X, Y]-$ коммутатор век- торных полей $X, Y$. Этот оператор также однозначно продолжается до дифференцирования $D(M)$, сохраняет тип тензоров и перестановочен со сверткой. В локальных координатах производная Ли тензора $T\left(T\left(\begin{array}{l}i_{1} \ldots i_{l} \\ j_{1} \ldots j_{m}\end{array}\right)\right.$ выражается так:

\end{enumerate}



\[

\begin{gathered}

L_{X} T=\xi^{k} \frac{\partial T_{j_{1} \ldots j_{m}}^{i_{1} \ldots i_{l}}}{\partial u^{k}}+T_{k \ldots j_{m}}^{i_{1} \ldots i_{l}} \frac{\partial \xi^{k}}{\partial u_{1}^{i_{1}}}+\ldots \\

\ldots-T_{j_{1} \ldots j_{m}}^{k \ldots i_{l}} \frac{\partial \xi^{i_{1}}}{\partial u^{k}}-\ldots

\end{gathered}

\]



\begin{enumerate}

  \setcounter{enumi}{2}

  \item Внешний дифференциал (внешняя производная) линейный оператор $d$, сопоставляющий внешней $\partial u ф-$ беренциальной борме (кососимметричному ковариантному тензору) степени $p$ форму такого же вида и степени $p+1$, удовлетворяющий условиям:

\end{enumerate}



\[

d\left(\omega_{1} \wedge \omega_{2}\right)=d \omega_{1} \wedge \omega_{2}+(-1)^{r} \omega_{1} \wedge d \omega_{2}, \quad d(d \omega)=0,

\]



где $\wedge$ - символ внешнего произведения, $r$ - степень $\omega_{1}$. В локальных координатах внешняя производная тензора $\omega\left(\omega_{i_{1} \ldots i_{p}}\right)$ выражается так:



\[

d \omega=\sum_{k=1}^{p+1}(-1)^{k} \frac{\partial \omega_{i_{1} \ldots \hat{i}_{k} \ldots i_{p+1}}^{\partial u_{k}}}{\partial{ }^{i_{k}}} .

\]



Оператор $d$ - обобщение оператора rot.



\begin{enumerate}

  \setcounter{enumi}{3}

  \item Кривизны тензор симметричного невырожденного дважды ковариантного тензора $g_{i j}$ представляет собой действие нек-рого нелинейного оператора $R$ :

\end{enumerate}



\[

g_{i j} \rightarrow R_{m l k}^{S}=\frac{\partial \Gamma_{k m}^{s}}{\partial u^{l}}-\frac{\partial \Gamma_{k l}^{s}}{\partial u^{m}}+\sum_{p}\left(\Gamma_{l p}^{s} \Gamma_{k m}^{p}-\Gamma_{m p}^{s} \Gamma_{k l}^{p}\right)

\]



где



\[

\Gamma_{j k}^{i}=\frac{1}{2} g i s\left(\frac{\partial g_{\jmath j s}}{\partial u^{k}}+\frac{\partial g_{k s}}{\partial u^{s}}-\frac{\partial g_{j k}}{\partial u^{s}}\right) .

\]



Лит.: [1] Р а ш е в с к и й П. К., Риманова геометрия и тензорный анализ, 3 изд., М., 1967; [2] С х о у т ен Я.-А., Тензорный анализ для физиков, пер. с англ., М., 1965; [3] М а кК о н н е л А.-Д., Введение в тензорный аналив, пер. с англ.,

М., 1963; [4] С к о л в н и к о в И. С., Тензорный анализ, пер. с англ., М., 1971 .



TEOPEMA - математическое утверждение, истинность к-рого установлена путем доказательства.



Понятие Т. развивалось п уточнялось вместе с понятием математич. доказательства. При использовании аксиоматического метода Т. рассматриваемой теории определяются как высказывания, выводимые чисто логич. путем из нек-рых заранее выбранных и фиксированных высказываний, называемых а к с и о м а м и. Поскольку аксиомы предполагаются истинными, то истинными должны быть и Т. Дальнейшее уточнение понятий доказательства и Т. связано с предпринятым в математич. логике исследованием понятия логического следствия, в результате чего для широкого класса математич. теорий процесс логич. вывода удалось свести к преобразованию формул, т. е. математич. утверждений, записанных на подходящем формализованном языке, по точно сформулированным правилам (вывода правилам), относящимся лишь к форме (а не к содержанию) предложений. В возникающих таким образом формальных теориях д о каз а т е ль с т в о м наз. конечная последовательность формул, каждая из к-рых либо является аксиомой, либо получается из нек-рых предыдущих формул этой последовательности по одному из правил вывода. Т. наз. формула, являющаяся последней формулой в нек-ром доказательстве.



Такос уточнение понятия $\mathrm{T}$. позволило получить, пользуясь строгими математич. методами, ряд важных





\end{document}